\section*{Capítulo \ref{ch:g10:stats} --- Estadística descriptiva}

\begin{solution}{exo:g10:stats:1}
Suma: $0+1+1+2+2+2+3+3+4+4+5+6+7+8+12 = 60$, luego la \omterm{def:g10:stats:mean}{media} es
$\frac{60}{15} = 4$. La serie ya está ordenada y $N = 15$ es impar: la
\omterm{def:g10:stats:median}{mediana} es el $8$.\textsuperscript{o} valor, $3$. \omterm{def:g10:stats:quartiles}{Recorrido}:
$12 - 0 = 12$.
\end{solution}

\begin{solution}{exo:g10:stats:2}
\omterm{def:g10:stats:series}{Frecuencias relativas}: $\frac{7}{50} = 0.14$, $\frac{9}{50} = 0.18$,
$\frac{8}{50} = 0.16$, $\frac{10}{50} = 0.20$, $\frac{6}{50} = 0.12$,
$\frac{10}{50} = 0.20$ (suman $1$). Resultado medio:
\[
\bar x = \frac{7 \times 1 + 9 \times 2 + 8 \times 3 + 10 \times 4
+ 6 \times 5 + 10 \times 6}{50}
= \frac{7 + 18 + 24 + 40 + 30 + 60}{50}
= \frac{179}{50} = 3.58 .
\]
\end{solution}

\begin{solution}{exo:g10:stats:3}
$N = 12$ (par): la \omterm{def:g10:stats:median}{mediana} es el \omterm{prop:g10:coordgeom:midpoint}{punto medio} del $6$.\textsuperscript{o} y
el $7$.\textsuperscript{o} valores, los dos iguales a $6$: \omterm{def:g10:stats:median}{mediana} $6$.

$Q_1$: $\frac{12}{4} = 3$, así que $Q_1$ es el $3$.\textsuperscript{er}
valor: $4$.

$Q_3$: $\frac{3 \times 12}{4} = 9$, así que $Q_3$ es el
$9$.\textsuperscript{o} valor: $8$.

Diagrama de caja: caja de $4$ a $8$ con la barra de la \omterm{def:g10:stats:median}{mediana} en $6$ y
bigotes de $1$ a $11$.
\end{solution}

\begin{solution}{exo:g10:stats:4}
Las $24$ notas suman $24 \times 11.5 = 276$. Con la nueva nota, el total
es $276 + 19 = 295$ para $25$ alumnos: nueva \omterm{def:g10:stats:mean}{media}
$\frac{295}{25} = 11.8$.
\end{solution}

\begin{solution}{exo:g10:stats:5}
Total actual: $5 \times 12 = 60$. Una \omterm{def:g10:stats:mean}{media} de $13$ en $6$ exámenes exige
un total de $78$, así que la $6$.\textsuperscript{a} nota debe ser
$78 - 60 = 18$. Una \omterm{def:g10:stats:mean}{media} de $14$ exigiría $84 - 60 = 24 > 20$: no es
alcanzable.
\end{solution}

\begin{solution}{exo:g10:stats:6}
\omterm{def:g10:stats:mean}{Media} mayor que $20$: por ejemplo, $1,\ 2,\ 3,\ 10,\ 40,\ 50,\ 60$, con
\omterm{def:g10:stats:median}{mediana} $10$ (el $4$.\textsuperscript{o} valor ordenado) y \omterm{def:g10:stats:mean}{media}
$\frac{166}{7} \approx 23.7$.

\omterm{def:g10:stats:mean}{Media} menor que $6$: por ejemplo, $0,\ 0,\ 0,\ 10,\ 10,\ 10,\ 10$, con
\omterm{def:g10:stats:median}{mediana} $10$ y \omterm{def:g10:stats:mean}{media} $\frac{40}{7} \approx 5.7$.

La \omterm{def:g10:stats:mean}{media} no puede bajar de $\frac{40}{7}$: para que la \omterm{def:g10:stats:median}{mediana} (el
$4$.\textsuperscript{o} valor ordenado) sea $10$, los valores de las
posiciones $4$ a $7$ deben ser todos $\geq 10$, así que el total es al
menos $40$ y la \omterm{def:g10:stats:mean}{media}, al menos $\frac{40}{7}$; el segundo ejemplo es el
extremo.

Estos ejemplos muestran que la \omterm{def:g10:stats:mean}{media} y la \omterm{def:g10:stats:median}{mediana} son en gran medida
independientes: conocer una dice poco de la otra, y por eso un buen
resumen informa de las dos.
\end{solution}

\begin{solution}{exo:g10:stats:7}
\emph{1.} \omterm{def:g10:stats:mean}{Media}:
$\frac{9 \times 1800 + 12000}{10} = \frac{16200 + 12000}{10} = 2820$.
\omterm{def:g10:stats:median}{Mediana}: ordenados, los $10$ salarios son nueve valores $1800$ y luego
$12000$; la \omterm{def:g10:stats:median}{mediana} es el \omterm{prop:g10:coordgeom:midpoint}{punto medio} del $5$.\textsuperscript{o} y el
$6$.\textsuperscript{o} valores, los dos $1800$: \omterm{def:g10:stats:median}{mediana} $1800$.

\emph{2.} La \omterm{def:g10:stats:median}{mediana} ($1800$) describe el salario típico: es lo que cobran
realmente $9$ de cada $10$ empleados. La \omterm{def:g10:stats:mean}{media} ($2820$) está tirada hacia
arriba por el único salario alto y no coincide con la nómina de nadie.
\end{solution}

\begin{solution}{exo:g10:stats:8}
Estatura total: $12 \times 168 + 18 \times 163 = 2016 + 2934 = 4950$ cm
para $30$ alumnos, así que la \omterm{def:g10:stats:mean}{media} es $\frac{4950}{30} = 165$ cm. Está
más cerca de la \omterm{def:g10:stats:mean}{media} de las chicas porque son más numerosas (\omterm{def:g10:stats:mean}{media}
ponderada, con pesos $12$ y $18$).
\end{solution}

\begin{solution}{exo:g10:stats:9}
\emph{1.} La nueva \omterm{def:g10:stats:mean}{media} es
\[
\bar y = \frac{y_1 + \dots + y_N}{N}
= \frac{(a x_1 + b) + \dots + (a x_N + b)}{N}
= \frac{a(x_1 + \dots + x_N) + Nb}{N}
= a \bar x + b .
\]

\emph{2.} Con $a = 1.8$ y $b = 32$: \omterm{def:g10:stats:mean}{media} $1.8 \times 20 + 32 = 68$ grados
Fahrenheit.
\end{solution}

\begin{solution}{pb:g10:stats:1}
\textbf{1.} \omterm{def:g10:stats:mean}{Media}: $\frac{2\,550}{7} \approx 364$ mil euros; \omterm{def:g10:stats:median}{mediana}: el
cuarto de los siete precios ordenados, $230$. La \omterm{def:g10:stats:median}{mediana} describe la
calle: seis de las siete casas no cuestan nada parecido a $364$.

\textbf{2.} Sin el caserón: \omterm{def:g10:stats:mean}{media} $\frac{1\,350}{6} = 225$, \omterm{def:g10:stats:median}{mediana}
$\frac{220 + 230}{2} = 225$. La \omterm{def:g10:stats:mean}{media} ha caído unos $139$; la \omterm{def:g10:stats:median}{mediana},
$5$. Un solo valor extremo puede arrastrar la \omterm{def:g10:stats:mean}{media} a donde sea; la
\omterm{def:g10:stats:median}{mediana} apenas se entera: es \emph{robusta}.

\textbf{3.} \omterm{def:g10:stats:mean}{Media} $= \frac{0 \times 20 + 1 \times 25 + 2 \times
30 + 3 \times 15 + 4 \times 10}{100} = \frac{170}{100} = 1.7$. Ninguna
familia tiene $1.7$ hijos, pero $1.7 \times 100 = 170$ \emph{sí} es el
número total exacto de hijos: la \omterm{def:g10:stats:mean}{media} codifica el total, repartido a
partes iguales; ese es su verdadero talento.

\textbf{4.} $N = 15$: \omterm{def:g10:stats:median}{mediana} $=$ $8$.\textsuperscript{o} valor $= 11$;
$Q_1 =$ $4$.\textsuperscript{o} valor $= 8$; $Q_3 =$
$12$.\textsuperscript{o} valor $= 14$.

\textbf{5.} \omterm{def:g10:stats:quartiles}{Recorrido} $19 - 4 = 15$: la dispersión completa, extremos
incluidos. \omterm{def:g10:stats:quartiles}{Recorrido intercuartílico} $14 - 8 = 6$: la anchura de la mitad
central, o sea, la dispersión sin el ruido de los extremos.

\textbf{6.} Un desplazamiento $+b$ mueve \emph{todos} los valores, así que
las diferencias entre valores (\omterm{def:g10:stats:quartiles}{recorrido}, \omterm{def:g10:stats:quartiles}{recorrido intercuartílico}) no
cambian: $(x_i + b) - (x_j + b) = x_i - x_j$. Un cambio de escala de razón
$a$ multiplica todas las diferencias por $a$: el \omterm{def:g10:stats:quartiles}{recorrido} y el \omterm{def:g10:stats:quartiles}{recorrido
intercuartílico} quedan multiplicados por $\abs a$.

\textbf{7.} Corrección aditiva: $6 \to 9$ y $18 \to 21$; el alumno flojo
gana un $50\,\%$ y el fuerte, un $17\,\%$, y el $21$ solo se sale
ligeramente de la escala de $20$ puntos. Multiplicativa: $6 \to 8$ y
$18 \to 24$; se conservan las razones, pero $24$ es imposible en una
escala de $20$: el problema del techo. Las correcciones aditivas favorecen
a los de abajo; las multiplicativas hacen estallar a los de arriba.

\textbf{8.} Total de notas: $20 \times 12 + 30 \times 10 = 540$, así que
la \omterm{def:g10:stats:mean}{media} conjunta es $\frac{540}{50} = 10.8$: los treinta alumnos de la
clase B pesan más que los veinte de la A; las \omterm{def:g10:stats:mean}{medias} se fusionan por \omterm{def:g10:stats:mean}{media}
\emph{ponderada}, nunca por \omterm{def:g10:stats:mean}{media} simple.

\textbf{9.} \omterm{def:g10:stats:median}{Mediana} de $A$: $2$; la de $B$: $5$. Lista fusionada
$\{0, 1, 2, 3, 10\}$: \omterm{def:g10:stats:median}{mediana} $2$. La \omterm{def:g10:stats:mean}{media} de las \omterm{def:g10:stats:median}{medianas} ponderada por
tamaños sería $\frac{3 \times 2 + 2 \times 5}{5} = 3.2 \neq 2$: las
\omterm{def:g10:stats:median}{medianas} no llevan ningún total, así que no existe fórmula de fusión; hay
que reordenar la lista entera.

\textbf{10.} \omterm{def:g10:stats:mean}{Media}: $\frac{31.3}{5} = 6.26$. \omterm{def:g10:stats:median}{Mediana}: $5.4$. \omterm{def:g10:stats:mean}{Media}
recortada: $\frac{5.3 + 5.4 + 5.5}{3} = 5.4$. Un juez entusiasta (o
corrupto) ha movido la \omterm{def:g10:stats:mean}{media} casi un punto; la \omterm{def:g10:stats:median}{mediana} y la \omterm{def:g10:stats:mean}{media}
recortada se han encogido de hombros: de ahí las reglas de recorte de los
deportes con jueces.

\textbf{11.} Fábrica A: la mitad central cabe en $0.2$ mm; fábrica B, en
$4$ mm. Misma \omterm{def:g10:stats:mean}{media}, fiabilidad opuesta: compra en A. Ha decidido el
\omterm{def:g10:stats:quartiles}{recorrido intercuartílico}; las \omterm{def:g10:stats:mean}{medias} solas comparan centros, nunca
regularidad.

\textbf{12.} Una \omterm{def:g10:stats:median}{mediana} muy por debajo de la \omterm{def:g10:stats:mean}{media} delata una
\emph{cola larga por la derecha}: casi todas las esperas cortas y unas
pocas catastróficas. Ejemplo: $\{5, 5, 10, 10, 95\}$, con \omterm{def:g10:stats:median}{mediana} $10$ y
\omterm{def:g10:stats:mean}{media} $\frac{125}{5} = 25$.

\textbf{13.} Significa que el $75\,\%$ de los bebés de esa edad son más
bajos (y el $25\,\%$, más altos): es un puesto, no una talla. ``Todos los
niños por encima de la \omterm{def:g10:stats:mean}{media}'' es imposible: si todos los valores
superaran a la \omterm{def:g10:stats:mean}{media}, su promedio superaría a la \omterm{def:g10:stats:mean}{media}, que es él mismo:
contradicción. Pero \emph{casi} todos sí pueden estarlo: un solo niño muy
bajo puede mantener la \omterm{def:g10:stats:mean}{media} por debajo de todos los demás. Por encima de
la \emph{\omterm{def:g10:stats:median}{mediana}}: nunca más de la mitad, por definición. A los puestos no
se los puede halagar; a los promedios, sí.

\textbf{14.} Nota típica: la misma, \omterm{def:g10:stats:median}{mediana} $11$ en las dos.
Homogeneidad de la mitad central: B es más compacta (\omterm{def:g10:stats:quartiles}{recorrido
intercuartílico} $3$ frente a $5$). Extremos: B es más salvaje (\omterm{def:g10:stats:quartiles}{recorrido}
$17$ frente a $13$). B es una clase de alumnos parecidos con unos pocos
valores atípicos; A está repartida de manera uniforme.

\textbf{15.} (1) ¿Es $5\,000$ la \omterm{def:g10:stats:mean}{media} o la \omterm{def:g10:stats:median}{mediana}, y me pueden dar la
\omterm{def:g10:stats:median}{mediana}? (2) ¿Cuál es la dispersión (los \omterm{def:g10:stats:quartiles}{cuartiles}), para saber hasta qué
punto lo ``típico'' es típico? (3) ¿A quién se ha contado: tamaño de la
muestra y forma de elegirla (la lección sobre encuestas del problema de
fin de semana del volumen anterior sobre cómo mienten los datos)?

\textbf{16.} Por ejemplo, $100, 150, 200, 250, 800$: \omterm{def:g10:stats:median}{mediana} $200$, \omterm{def:g10:stats:mean}{media}
$\frac{1\,500}{5} = 300$.

\textbf{17.} $\sum (x_i - \bar x) = \sum x_i - N \bar x = N\bar x -
N\bar x = 0$. Comprobación: $-200 - 150 - 100 - 50 + 500 = 0$. La \omterm{def:g10:stats:mean}{media} es
el punto donde se equilibran las desviaciones: el centro de gravedad de
los datos (el $G$ del problema de fin de semana del volumen anterior sobre
el baricentro, en una dimensión).

\textbf{18.} En la \omterm{def:g10:stats:median}{mediana} $4$: $3 + 1 + 0 + 5 + 9 = 18$ km. En la \omterm{def:g10:stats:mean}{media}
$6$: $5 + 3 + 2 + 3 + 7 = 20$ km. En $5$: $4 + 2 + 1 + 4 + 8 = 19$ km.
Gana la \omterm{def:g10:stats:median}{mediana}: alejarse de ella siempre deja atrás más tiendas de las
que acerca, así que el total solo puede crecer.

\textbf{19.} $0.3^{10} \approx 6 \times 10^{-6}$: unos seis pilotos
\emph{por millón}. Entre $4\,000$ pilotos, el número esperado es $0.02$,
es decir, cero, tal como encontró el teniente Daniels. La Fuerza Aérea
dejó de diseñar para el hombre medio y lo hizo todo ajustable: asientos,
pedales, correas; diseñar para la \emph{dispersión}, no para el centro.

\textbf{20.} Totales: la \emph{\omterm{def:g10:stats:mean}{media}}. Individuo típico: la
\emph{\omterm{def:g10:stats:median}{mediana}}. Equidad de la dispersión: los \emph{\omterm{def:g10:stats:quartiles}{cuartiles}} (y el
\omterm{def:g10:stats:quartiles}{recorrido intercuartílico}). Contaminación: la \emph{\omterm{def:g10:stats:mean}{media} recortada}. Y la
frase de cierre: la persona \omterm{def:g10:stats:mean}{media} no existe, pero la persona \omterm{def:g10:stats:median}{mediana} casi
sí.
\end{solution}
